\begin{thebibliography}{30}\setlength{\itemsep}{0pt}
\bibitem{pib23} Press Information Bureau, Ministry of Heavy Industries, ``6,586 operational public EV charging stations in India,'' Press Release, New Delhi, Mar. 24, 2023 (Annexure A). [Online]. Available: \url{https://www.pib.gov.in/PressReleaseIframePage.aspx?PRID=1910392&reg=48&lang=2}

\bibitem{pib24} Press Information Bureau, Ministry of Heavy Industries, ``12,146 public EV charging stations operational across the country,'' Press Release, New Delhi, Feb. 6, 2024 (Annexure II). [Online]. Available: \url{https://www.pib.gov.in/PressReleaseIframePage.aspx?PRID=2003003&reg=48&lang=2}

\bibitem{pib25} Press Information Bureau, Ministry of Heavy Industries, ``Electric vehicle charging stations,'' Press Release, New Delhi, Dec. 16, 2025. [Online]. Available: \url{https://www.pib.gov.in/PressReleasePage.aspx?PRID=2204635&reg=3&lang=1}

\bibitem{sansad} Ministry of Heavy Industries, ``Installation of electric vehicle charging stations,'' Rajya Sabha Unstarred Question No. 2333, Parliament of India, answered Aug. 7, 2026. [Online]. Available: \url{https://sansad.in/getFile/annex/271/AU2333_V2LOXS.pdf?source=pqars}

\bibitem{emobility} eMobilityPlus, ``West Bengal e-rickshaw drive uncovers hidden EV fleet in India,'' Feb. 9, 2026. [Online]. Available: \url{https://emobilityplus.com/2026/02/09/west-bengal-e-rickshaw-drive-uncovers-hidden-ev-fleet-in-india}

\bibitem{vahan} Ministry of Road Transport and Highways (MoRTH), ``Vahan public dashboard: vehicle registrations,'' Government of India. [Online]. Available: \url{https://analytics.parivahan.gov.in/analytics/publicdashboard/vahan}. Accessed: Oct. 9, 2026.

\bibitem{cpcb} Central Pollution Control Board (CPCB), ``Central Control Room for Air Quality Management: CAAQMS dashboard,'' Government of India. [Online]. Available: \url{https://airquality.cpcb.gov.in/ccr/#/}. Accessed: Oct. 9, 2026.

\bibitem{ppac} Petroleum Planning and Analysis Cell (PPAC), ``Historical retail selling prices of petrol and diesel,'' Ministry of Petroleum and Natural Gas, Government of India. [Online]. Available: \url{https://ppac.gov.in}. Accessed: 2026.

\bibitem{nasa} NASA Langley Research Center, ``Prediction Of Worldwide Energy Resources (POWER) Project: Temporal API (point).'' [Online]. Available: \url{https://power.larc.nasa.gov/}. Accessed: Nov. 9, 2026.

\bibitem{mospi} Ministry of Statistics and Programme Implementation (MoSPI), ``National accounts statistics: State domestic product'' [Data set], e-Sankhyiki portal, Apr. 17, 2026. [Online]. Available: \url{https://esankhyiki.mospi.gov.in}

\bibitem{sierzchula} W. Sierzchula, S. Bakker, K. Maat, and B. van Wee, ``The influence of financial incentives and other socio-economic factors on electric vehicle adoption,'' \emph{Energy Policy}, vol. 68, pp. 183--194, 2014, doi: 10.1016/j.enpol.2014.01.043.

\bibitem{li} S. Li, L. Tong, J. Xing, and Y. Zhou, ``The market for electric vehicles: Indirect network effects and policy design,'' \emph{J. Assoc. Environ. Resour. Econ.}, vol. 4, no. 1, pp. 89--133, 2017, doi: 10.1086/689702.

\bibitem{devarasan} E. Devarasan, D. Nagarajan, and J. Rachel, ``Advancing sustainable mobility in India with electric vehicles: Market trends and machine learning insights,'' \emph{Front. Energy Res.}, vol. 13, Art. no. 1500515, 2025, doi: 10.3389/fenrg.2025.1500515.

\bibitem{rohit} K. Rohit and A. Verma, ``Forecasting electric two-wheeler demand in the Indian supply chain using machine learning models,'' \emph{Proc. Inst. Mech. Eng. D, J. Automobile Eng.}, vol. 240, no. 8, pp. 5076--5087, 2026, doi: 10.1177/09544070251372178.

\bibitem{radhakrishnan} H. S. Radhakrishnan, A. Raj, and P. Bansal, ``Wheels of change: Forecasting regional electric two-wheeler adoption in India,'' \emph{Transp. Res. D, Transp. Environ.}, vol. 159, Art. no. 105504, 2026.

\bibitem{afandizadeh} S. Afandizadeh, D. Sharifi, N. Kalantari, and H. Mirzahossein, ``Using machine learning methods to predict electric vehicles penetration in the automotive market,'' \emph{Sci. Rep.}, vol. 13, Art. no. 8345, 2023, doi: 10.1038/s41598-023-35366-3.

\bibitem{chen} T. Chen and C. Guestrin, ``XGBoost: A scalable tree boosting system,'' in \emph{Proc. 22nd ACM SIGKDD Int. Conf. Knowl. Discovery Data Mining}, 2016, pp. 785--794.

\bibitem{ke} G. Ke \emph{et al.}, ``LightGBM: A highly efficient gradient boosting decision tree,'' in \emph{Adv. Neural Inf. Process. Syst. 30}, 2017, pp. 3146--3154.

\bibitem{grinsztajn} L. Grinsztajn, E. Oyallon, and G. Varoquaux, ``Why do tree-based models still outperform deep learning on typical tabular data?'' in \emph{Adv. Neural Inf. Process. Syst. 35 (Datasets and Benchmarks Track)}, 2022.

\bibitem{lundberg} S. M. Lundberg and S.-I. Lee, ``A unified approach to interpreting model predictions,'' in \emph{Adv. Neural Inf. Process. Syst. 30}, 2017, pp. 4765--4774.

\bibitem{makridakis} S. Makridakis, E. Spiliotis, and V. Assimakopoulos, ``Statistical and machine learning forecasting methods: Concerns and ways forward,'' \emph{PLoS ONE}, vol. 13, no. 3, Art. no. e0194889, 2018, doi: 10.1371/journal.pone.0194889.

\bibitem{neweywest} W. K. Newey and K. D. West, ``A simple, positive semi-definite, heteroskedasticity and autocorrelation consistent covariance matrix,'' \emph{Econometrica}, vol. 55, no. 3, pp. 703--708, 1987.

\bibitem{hyndman} R. J. Hyndman and G. Athanasopoulos, \emph{Forecasting: Principles and Practice}, 3rd ed. Melbourne, Australia: OTexts, 2021. [Online]. Available: \url{https://otexts.com/fpp3/}

\end{thebibliography}
